% ==================================================================
% id:       cond-008
% topic:    הסתברות מותנית ואי-תלות
% subtopic: הסתברות מותנית
% solution: full
% ==================================================================
בגן ילדים מגישים בכל בוקר שוקו ועוגיות. 45\% מהילדים בוחרים שוקו, 60\% בוחרים עוגייה, ו-30\% בוחרים גם שוקו וגם עוגייה. בוחרים באקראי ילד או ילדה מהגן, ומתברר שבחרו עוגייה. מה ההסתברות שלא בחרו שוקו?
\begin{solbox}
נסמן $H$ = ״בחר/ה שוקו״, $C$ = ״בחר/ה עוגייה״. נתון $\PP(C)=0.6$ ו-$\PP(H\cap C)=0.3$, ומחפשים את $\PP(H^C\mid C)$.

המאורע $C$ הוא האיחוד הזר של $H\cap C$ ושל $H^C\cap C$, ולכן $\PP(H^C\cap C)=0.6-0.3=0.3$. לפי ההגדרה של הסתברות מותנית
\[
\PP(H^C\mid C)=\frac{\PP(H^C\cap C)}{\PP(C)}=\frac{0.3}{0.6}=\ans{0.5}.
\]
שימו לב: הנתון ש-45\% בוחרים שוקו לא נחוץ.
\end{solbox}
